% ECONOMICS_PROBLEM: {"id": "D12-516", "title": "Privacy valuation", "jel_code": "D12", "source_id": "OP-0406"}
% !TeX program = xelatex
\documentclass[11pt,a4paper]{article}
\usepackage[margin=23mm]{geometry}
\usepackage{fontspec}
\setmainfont{Latin Modern Roman}
\usepackage{amsmath,amssymb,mathtools}
\usepackage{xcolor,enumitem,booktabs,longtable,array,xurl,hyperref}
\definecolor{muted}{HTML}{53636C}
\hypersetup{colorlinks=true,urlcolor=blue,linkcolor=blue}
\setlength{\parindent}{0pt}
\setlength{\parskip}{5pt}
\newcommand{\R}{\mathbb R}
\newcommand{\E}{\mathbb E}
\newcommand{\N}{\mathbb N}
\newcommand{\ind}{\mathbf 1}
\newcommand{\Var}{\operatorname{Var}}
\newcommand{\Cov}{\operatorname{Cov}}
\newcommand{\supp}{\operatorname{supp}}
\DeclareMathOperator*{\argmin}{arg\,min}
\DeclareMathOperator*{\argmax}{arg\,max}
\newcommand{\entrymeta}[3]{{\sffamily\small\color{muted}\textbf{\texttt{#1}}\quad|\quad #2\par #3\par}}
\newcommand{\Problemref}[1]{\texttt{#1}}
\newcommand{\JELref}[2]{\texttt{#2} (JEL \texttt{#1})}
\begin{document}
\section*{Privacy valuation}
\textbf{Problem ID:} \texttt{D12-516}\quad \textbf{JEL:} \texttt{D12}\par
% BEGIN ECONOMICS STATEMENT
\label{problem:OP-0406}
{\sffamily\small\color{muted}Original subject group: Digital, platform and data economics\par}
\label{definitions:problem:30007688}
\textbf{Audit classification:} Background; proposed extension. The publication exists and is relevant background. No exact published open theorem for this CV identification exercise was verified.
\subsection*{Definitions and precise research target}
Consider adult users of an online service in a fixed country and enrollment period. A data-sharing contract \(d\in\{0,1\}\) either limits collection to essential service data or permits a specified additional use; the contract's actual data flows are auditable. Let \(u_i(c,d,\ell)\) be user \(i\)'s utility from baseline consumption \(c>0\), sharing regime \(d\), and future loss \(\ell\), with utility strictly increasing in \(c\). The additional use induces an uncertain loss distribution \(F_i\) over a fixed five-year horizon, including discrimination and unwanted disclosure. Define compensating variation \(CV_i\) by \(E_{F_i}[u_i(c+CV_i,1,\ell)]=u_i(c,0,0)\), with expectation over the user's stated informed loss beliefs. Require finite expected utility and \(c+CV_i>0\); report nonexistence if no finite compensation attains the comparison utility. The target is the population distribution of \(CV_i\), alongside sensitivity to alternative credible loss distributions. Incentivized contract choices, comprehension checks and belief elicitation provide distinct observations; default effects cannot automatically be interpreted as valuations. Determine which additional information treatments, repeated choices or structural restrictions distinguish preferences from misunderstanding and friction. When distant harms cannot be learned from the observation window, bound valuations under explicit probability and loss limits. This is a proposed identification exercise rather than a verified universal unresolved privacy-valuation theorem.

\textbf{Definitions, assumptions and mathematical qualifications.} The contract defines actual additional data recipients, retention, permitted inference and the five-year loss variable, rather than a vague sharing label. Normalize \(d=0\) loss to zero only if baseline risks are incorporated in baseline utility. Define \(f_i(v)=\int u_i(c+v,1,\ell)dF_i(\ell)\) on \(v>-c\). Continuity, strict increase and integrability imply at most one \(CV_i\); existence further requires \(u_i(c,0,0)\) to lie in the range of \(f_i\). Compensation may be negative when extra services outweigh privacy losses. The population CV distribution is \(G(v)=\Pr(CV_i\leq v)\), with nonexistence separately tabulated. Stated subjective loss beliefs and externally estimated loss frequencies are distinct estimands.
\subsection*{Variants and assumption changes}
\begin{enumerate}
\item \textbf{Editorial, informed benchmark.} Under quasi-linear utility \(u_i(c,d,\ell)=c+b_i(d)-h_i(d)-\ell\), \(CV_i=b_i(0)-b_i(1)+h_i(1)-h_i(0)+E_{F_i}\ell\); the issue is measuring these components.
\item \textbf{Editorial, uncertain harms.} Let loss laws belong to a stated set with bounded probability and monetary severity; combine randomized information and comprehension measures to bound informed valuation without treating default behavior as direct willingness to pay.
\end{enumerate}
\subsection*{References and source audit}
\textbf{Checked source.} NBER Working Paper 30943, February 2023, indexed primary abstract. The publication supports privacy cost--benefit and externality background. Primary metadata and abstract indexed; direct PDF returned HTTP 403. An exact open-question passage was not verified. \url{https://www.nber.org/papers/w30943}.
\begin{enumerate}\raggedright
\item\hypertarget{definitions:source:30007688:1}{} Avi Goldfarb; Verina F. Que. 2023. \emph{The Economics of Digital Privacy}. NBER Working Paper 30943, February 2023, indexed primary abstract.
\url{https://www.nber.org/papers/w30943}.
DOI: \url{https://doi.org/10.3386/w30943}.
\end{enumerate}

\subsection*{Common conventions: Source mathematical conventions}

These conventions apply to each entry unless explicitly replaced.

\textbf{Models and identification.} A primitive model \(m\) specifies all probability laws, feasible choices, information, equilibrium and measurement rules used by its target. The admissible class \(\mathcal M\) is fixed independently of outcomes. If \(P_O(m)\) is its observable probability law and \(\tau(m)\) the target, its identified set at observable law \(P\) is
\[
 \mathcal I_\tau(P)=\{\tau(m):m\in\mathcal M,\ P_O(m)=P\}.
\]
Point identification means that this set is a singleton. An empty set signals incompatibility of data and assumptions. Coordinate bounds need not describe the joint set. Bounds are sharp only if every retained target is attainable or arbitrarily approachable by compatible models. Finite bounds require explicit restrictions on unobserved quantities; unrestricted confounding, hidden wealth, extrapolation or arbitrary equilibrium selection can yield uninformative sets.

\textbf{Causal interventions.} A potential outcome \(Y_i(p)\) is the outcome under a complete policy regime \(p\), including timing, financing, information and market spillovers. The target population and units are fixed before treatment. With interference, use the complete assignment vector or a justified exposure mapping; individual no-interference notation alone cannot identify an economy-wide intervention. A difference of mean potential outcomes does not require the cross-world joint distribution. Conditioning on post-treatment migration, survival, enrollment or employment changes the estimand and needs separate justification.

\textbf{Statistical statements.} An estimator, confidence set, forecast or test specifies a sampling law, model class and loss/coverage criterion. Uniform \((1-\alpha)\)-coverage means \(\inf_{m\in\mathcal M}\Pr_m\{\tau(m)\in C_n\}\geq1-\alpha\), or an explicitly asymptotic version. The whole parameter space trivially has coverage, so informativeness requires a stated width, risk or power objective. A forecasting improvement is relative to a named benchmark, information set and evaluation population.

\textbf{Equilibrium and optimization.} An equilibrium comprises feasible plans, optimality, beliefs where needed, budget constraints and all market/resource clearing. The primitives or a stated benchmark define each correspondence \(\mathcal E(p;m)\). Nonemptiness is assumed or proved, never inferred just from compact policy sets. A selected-equilibrium target uses a declared selection rule; a robust target quantifies over all permitted equilibria. Use suprema or infima unless attainment is established. An \(\varepsilon\)-optimal policy is feasible and within \(\varepsilon\) of the finite optimum value. Report an empty feasible set separately.

\textbf{Welfare and accounting.} Normative utility, interpersonal normalization, social weights, numeraire, discounting and the use of public receipts are primitives. Behavioral utility need not coincide with the welfare criterion. Household welfare includes ownership income and rents once; adding profits again to compensating variation generally double-counts them. A monetary equivalent is defined by an explicit transfer and price convention, with monotonicity and range conditions. Average effects, mechanism contributions, transfers and welfare losses are different targets. Infinite sums require convergence and dynamic feasible plans require terminal or no-Ponzi conditions.

\textbf{Source versions.} A working-paper date, revision date and journal publication year are distinguished. A DOI for a journal version is not automatically the DOI of an earlier manuscript. Locators refer to the stated version and printed pages where specified. A metadata-only check does not verify a claim about a full-text conclusion. Every entry's source subsection narrows the provenance claim accordingly.
% END ECONOMICS STATEMENT
\end{document}
